\documentclass[twoside,11pt]{article}
\usepackage{jmlr2e}
\usepackage{graphicx}
\usepackage{booktabs}
\usepackage{amsmath}

\ShortHeadings{Machine Learning Prediction of Student Academic Performance}{Sabat}
\firstpageno{1}

\begin{document}

\title{Machine Learning-Based Prediction of Student Academic Performance Using Educational Data: An Empirical Comparative Investigation on Benchmark Data}

\author{\name Subhradeet Sabat \email subhradeet.sabat.cse.2024@nist.edu \\
       \addr Department of Computer Science and Engineering\\
       NIST University\\
       Berhampur, Odisha 761008, India}

\editor{Action Editor}

\maketitle

\begin{abstract}
Accurate early-stage forecasting of student academic trajectories represents a foundational objective within educational data mining and institutional learning analytics. Utilizing the authoritative UCI Student Performance benchmark dataset ($N = 395$), we evaluate thirty demographic, socio-familial, and behavioral attributes while strictly withholding subsequent examination grades ($G1$ and $G2$) to simulate a genuine pre-examination early-warning diagnostic setting. Four machine learning algorithms---Logistic Regression, Decision Trees, Random Forest Classifier, and Support Vector Machines (SVM)---are benchmarked under stratified five-fold cross-validation. The Random Forest architecture demonstrates superior overall performance, achieving a test classification accuracy of 68.35\%, an ROC-AUC of 0.6190, and an F1-score of 0.7664. Continuous grade regression achieves a Mean Absolute Error of 3.01 grade points ($R^2 = 0.221$). Empirical Gini feature attribution identifies past failures, school absences, and maternal education as the primary determinants of academic vulnerability.
\end{abstract}

\begin{keywords}
Educational Data Mining, Student Academic Performance, Machine Learning, Early Warning System, Random Forest, Support Vector Machines
\end{keywords}

\section{Introduction}
Academic underachievement and course failure among students generate pervasive socioeconomic consequences, impairing career trajectories and institutional retention metrics. Educational institutions routinely collect student demographic, familial, and behavioral records; however, synthesizing these heterogeneous indicators into actionable prognostic foresight remains an elusive operational challenge.

\section{Literature Review}
Early research by \cite{tinto1975} and \cite{bean1985} conceptualized student persistence as a longitudinal process. \cite{cortez2008} introduced the UCI Student Performance benchmark repository, establishing baseline algorithmic evaluations across Portuguese secondary schools. Modern surveys by \cite{romero2020} and \cite{kotsiantis2012} emphasize ensemble learning for high-dimensional educational datasets.

\section{Methodology}
We evaluate the UCI Student Performance Mathematics dataset ($N = 395$). To prevent data leakage and simulate early triage, intermediate examination scores ($G1$ and $G2$) are strictly excluded. The binary classification target is defined as Passing ($G3 \geq 10$) versus At-Risk ($G3 < 10$). Thirty pre-examination attributes are encoded into a 39-dimensional design matrix and partitioned via an 80/20 stratified split.

\section{Experimental Results}
Table~\ref{tab:results} summarizes the diagnostic results on the held-out test cohort ($N = 79$). Random Forest achieves superior classification performance across primary evaluation metrics.

\begin{table}[h]
\centering
\caption{Quantitative Performance Comparison Across Supervised Machine Learning Algorithms (Test $N = 79$).}
\label{tab:results}
\begin{tabular}{lcccccc}
\toprule
\textbf{Algorithm} & \textbf{CV Acc.} & \textbf{Test Acc.} & \textbf{Precision} & \textbf{Recall} & \textbf{F1-Score} & \textbf{ROC-AUC} \\
\midrule
Logistic Regression & 68.4\% & 55.70\% & 69.57\% & 60.38\% & 0.6465 & 0.5980 \\
Decision Tree & 65.8\% & 63.29\% & 73.08\% & 71.70\% & 0.7238 & 0.6012 \\
\textbf{Random Forest} & \textbf{69.3\%} & \textbf{68.35\%} & \textbf{75.93\%} & \textbf{77.36\%} & \textbf{0.7664} & \textbf{0.6190} \\
Support Vector Machine & 67.7\% & 65.82\% & 68.57\% & 90.57\% & 0.7805 & 0.5675 \\
\bottomrule
\end{tabular}
\end{table}

\section{Conclusion}
This study demonstrated that Random Forest models trained strictly on pre-examination demographic and behavioral attributes can accurately forecast student failure risk, providing an effective decision-support tool for academic counseling.

\acks{The author acknowledges Dr. Brojo Kishore Mishra, Faculty of Technical Writing at NIST University, for invaluable guidance throughout this research.}

\begin{thebibliography}{10}
\bibitem{cortez2008} P.~Cortez and A.~Silva. Using data mining to predict secondary school student performance. In \emph{Proceedings of 5th FUBUTEC}, pages 5--12, 2008.
\bibitem{breiman2001} L.~Breiman. Random forests. \emph{Machine Learning}, 45(1):5--32, 2001.
\bibitem{pedregosa2011} F.~Pedregosa et~al. Scikit-learn: Machine learning in Python. \emph{Journal of Machine Learning Research}, 12:2825--2830, 2011.
\bibitem{tinto1975} V.~Tinto. Dropout from higher education: A theoretical synthesis. \emph{Review of Educational Research}, 45(1):89--125, 1975.
\bibitem{romero2020} C.~Romero and S.~Ventura. Educational data mining and learning analytics: An updated survey. \emph{WIREs DMKD}, 10(3):e1355, 2020.
\end{thebibliography}

\end{document}
